\section{The chord lemma: inventory and checks}
\label{app:A}

\emph{The proof} of Lemma~\ref{lem:chord} is in Sec.~\ref{sec:sandwich}; here we record what it is applied to.

\emph{The inventory.}
Table~\ref{tab:inventory} lists, for each line of~\cite{I}, the basic functions of $\eta$ that it contains.
Each loss and target line is a sum of products of the six functions of Lemma~\ref{lem:chord} (and of $1+\eta$) with non-negative $\eta$-free coefficients; each gain is $\eta$-free or carries the sandwich factor fl (and, for P2z, its remainder ringf).
The table is not written by hand: a parser reads the source of the lines and types every sub-expression by the grammar of Sec.~\ref{sec:sandwich}, and the table is its output \status{computer-assisted}.
\begin{itemize}
\item The types: $\eta$-free and non-negative (numbers and the state quantities $r,B,\kappa,P,\mu,b,k_E,g_s$, which do not involve $\eta$); \emph{convex} (non-negative, non-decreasing, convex in $\eta$): $\eta$, sums and products of convex and $\eta$-free terms, powers $\ge1$, division by $\eta$-free positive quantities or by powers of $1-\eta$, and $\sqrt{\mathrm{CZ}/c}$; the sandwich factor $\mathrm{fl}=1/(1+\eta)$; $\mathrm{ringf}=1-(\text{convex})$; anything else is rejected ($1-\eta$ in a numerator, a convex term in a denominator, the square root of another convex term).
\item The checks: every loss and target line is convex, every gain is $\eta$-free, $\mathrm{fl}\times$($\eta$-free) or $\mathrm{ringf}\cdot\mathrm{fl}\times$($\eta$-free); the list agrees with a hand-written one. The first hand-written list had three wrong rows (N4G does not contain $\eta$ itself; B2s and uB2s contain $\sqrt{\mathrm{CZ}}$ and not $\eta^2/(1-\eta)^2$); they were corrected.
\item Power: a loss line multiplied by $1-\eta$ (N2), $1+\eta$ replaced by $1/(1+\eta)$ (A1w), and $\sqrt\eta$ in d3, are rejected.
\item Numbers: the $83$ products of up to three of the six functions, on $201$ points of $\lambda\in[\tfrac45,1]$, lie under their chords (smallest difference $-9\times10^{-16}$, rounding); the loss lines of~\cite{I} at $60$ random states for each of $\ell=55.26$, $960.8$, $10^6$ and five values of $\lambda$ ($27\,105$ comparisons) lie under their chords (smallest relative difference $-7\times10^{-64}$, rounding of the decimal arithmetic).
\end{itemize}

\begin{table}
\caption{The inventory of Lemma~\ref{lem:chord}: the basic functions of $\eta$ in each of the 37 lines of~\cite{I}, read from the code by the parser (Sec.~\ref{sec:sandwich}). ``---'': the line does not contain $\eta$. Gains: fl $=1/(1+\eta)$ and ringf $=1-q$ are the sandwich and its remainder, split as in Sec.~\ref{sec:sandwich}. The line B1m is that of this paper, with $X_{1,\rm B1}(\eta)=U_1/(1-\eta)+\sqrt{\rm CZ}$ in place of $X_1$, read by the parser from the code transformed in that one place.}
\label{tab:inventory}
\begin{ruledtabular}
\small
\begin{tabular}{lll}
line & kind & basic functions\\ \hline
P1 & gain & ---\\
P2z & gain & fl, ringf: $\frac{1+\eta}{1-\eta}$\\
P3z & gain & ---\\
P4z & gain & ---\\
hFz & gain & fl\\
cKz & gain & fl\\
d1 & target & $\eta$\\
d2 & target & $\frac1{1-\eta}$\\
d3 & target & $\eta$\\
d4 & target & $\frac1{1-\eta}$\\
N1 & loss & $\eta$\\
N2 & loss & $\eta$\\
N3 & loss & $\frac1{1-\eta}$, $\eta$\\
N4 & loss & $\frac1{(1-\eta)^2}$, $\eta$, $\eta^2$\\
N4G & loss & $\frac1{1-\eta}$, $\frac1{(1-\eta)^2}$, $\eta^2$\\
ring2 & loss & $\frac1{1-\eta}$, $\eta^2$\\
M0Xi & loss & ---\\
A1v & loss & ---\\
A1w & loss & $\eta$\\
A2k & loss & $\frac1{1-\eta}$\\
A2a & loss & $\eta$\\
A2b & loss & $\frac1{1-\eta}$, $\eta^2$\\
B1w1 & loss & ---\\
B1w2 & loss & $\frac1{1-\eta}$, $\eta^2$\\
B2a & loss & $\frac1{(1-\eta)^2}$\\
B2b & loss & $\frac1{(1-\eta)^2}$, $\eta^2$\\
pair & loss & $\frac1{1-\eta}$\\
hneg1 & loss & $\eta$\\
lossK & loss & $\frac1{1-\eta}$\\
B2c & loss & $\frac1{1-\eta}$\\
B2s & loss & $\frac1{1-\eta}$, $\sqrt{\rm CZ}$\\
B1m & loss & $\frac1{1-\eta}$, $\sqrt{\rm CZ}$ (from $X_{1,\rm B1}$)\\
B1c & loss & $\frac1{1-\eta}$, $\frac1{(1-\eta)^2}$, $\eta^2$, $\sqrt{\rm CZ}$\\
B1l & loss & $\frac1{1-\eta}$\\
uB2c & loss & $\frac1{1-\eta}$\\
uB2s & loss & $\frac1{1-\eta}$, $\sqrt{\rm CZ}$\\
FzK & loss & $\frac1{1-\eta}$\\
\end{tabular}
\end{ruledtabular}
\end{table}


\emph{The split where the sandwich meets a remainder.}
P2z is $\mathrm{ringf}\cdot\mathrm{fl}\cdot G/(\mu(2+r)4r)$ with $\mathrm{ringf}=1-q$, $q=K^2(1+\eta)/(8(1-\eta)B)$.
With fl replaced by $\lambda$: $\lambda(1-q)\ge\lambda-q$, and $q$ is a remainder of the chord lemma (it contains $(1+\eta)/(1-\eta)$).
The same split is used in $\theta_0=\sqrt{(1-q/\lambda)(1-z_2/2)}$ of step (m1), which is the ratio of the gain $\lambda-q$ to the target $\lambda$.
d1 and d3 are split as in Sec.~\ref{sec:sandwich}.
